\section{PROBLEMA}

A resolução de falhas em ambientes de Infraestrutura como Código é um processo que demanda tempo significativo e conhecimento especializado. Quando o Terraform apresenta erros durante o provisionamento, os engenheiros precisam: (1) interpretar o log de erro, que pode ser extenso e conter múltiplas camadas de informação; (2) correlacionar a mensagem de erro com o trecho específico do código \texttt{.tf} que causou a falha; (3) consultar a documentação oficial para identificar a correção adequada; e (4) validar que a correção proposta não introduz novos problemas.

Este processo, quando realizado manualmente, é propenso a erros e altamente dependente da experiência individual do engenheiro. Em equipes com profissionais menos experientes, a resolução de falhas simples pode consumir horas, enquanto engenheiros sêniores frequentemente identificariam a causa raiz em minutos. Essa disparidade cria um gargalo operacional e evidencia a necessidade de ferramentas que democratizem o acesso ao conhecimento especializado.

Adicionalmente, soluções baseadas em APIs comerciais de IA (como OpenAI ou Anthropic) apresentam riscos de segurança ao exigir o envio de código de infraestrutura e logs operacionais para servidores externos, o que pode violar políticas de segurança corporativa e regulamentações de proteção de dados.

Diante desse cenário, o problema de pesquisa formula-se como: \textbf{``Como modelos de Inteligência Artificial Generativa, executados localmente, podem ser utilizados para automatizar a análise de logs de falhas em código Terraform, identificando a causa raiz e sugerindo correções aderentes às melhores práticas, sem comprometer a segurança dos dados da infraestrutura?''}
